% GENERATED FILE. Edit canonical lessons, then run npm run generate:books.
% canonical-source-hash: fnv1a64:b79920fc7d4aa60a
% canonical-lessons: TA-C554-atikapatcam, TA-C554-mutintavarai, TA-C554-malaiyerram, TA-C554-minpiti, TA-C554-pommalattam

\chapter{Weekend Plans}
\label{ch:ta-weekend-plans}
\hlchaptermodality{554}

\begin{chapteropening}
\textbf{By the end of this chapter:} \emph{I can talk about weekend plans and what we will do if the weather allows.}

Complete the last lesson of chapter 554: I can talk about weekend plans and what we will do if the weather allows.
\end{chapteropening}

\section[\texorpdfstring{\ta{அதிகபட்சம்} (atikapaṭcam)}{அதிகபட்சம் (atikapaṭcam)}]{\ta{அதிகபட்சம்} (atikapaṭcam) — at most; a maximum}

\label{lesson:TA-C554-atikapatcam}



\begin{quote}
Before the new one: say the Tamil for an alternative, then the Tamil for approximately.
\end{quote}

\subsection*{You'll want to know: \ta{அதிகபட்சம்}}
\textbf{\ta{அதிகபட்சம்}} — \emph{atikapaṭcam} — \textquotedblleft{}at most; a maximum\textquotedblright{}.

\begin{quote}
\ta{போக்குவரத்தைப்} \ta{பொறுத்து} \ta{அதிகபட்சம்} \ta{இரண்டு} \ta{மணி} \ta{நேரத்தில்} \ta{திரும்பி} \ta{வருவோம்}.
\end{quote}

\begin{quote}
\emph{pōkkuvarattaip poṟuttu atikapaṭcam iraṇṭu maṇi nērattil tirumpi varuvōm.}
\end{quote}

\begin{quote}
Depending on the traffic, we will be back in two hours at most.
\end{quote}

The pair of \ta{குறைந்தபட்சம்}, at least.

\begin{cousinweb}
\textbf{\ta{அதிகபட்சம்}} (\emph{atikapaṭcam}) is built from Sanskrit \emph{adhika}, \textquotedblleft{}additional, more\textquotedblright{} (CDIAL 250), and \emph{pakṣa}, \textquotedblleft{}a side, a wing\textquotedblright{} (CDIAL 7627): the side of more.
\end{cousinweb}

\begin{grammarlens}[title={What you've built}]
One more word for saying what happens if.
\end{grammarlens}

\subsection*{Your turn}
\begin{itemize}
\raggedright
  \item \emph{Say it:} \emph{atikapaṭcam}
  \item \emph{Say it:} \emph{atikapaṭcam}, once more
  \item \emph{Say it:} \emph{pōkkuvarattaip poṟuttu atikapaṭcam iraṇṭu maṇi nērattil tirumpi varuvōm.} — Depending on the traffic, we will be back in two hours at most.
  \item \emph{Recall:} say the Tamil for an alternative, then the Tamil for approximately, then say \emph{atikapaṭcam} again
\end{itemize}

\begin{tcolorbox}[breakable,colback=teal!4,colframe=teal!35!black,title={Before you move on}]
What does \ta{அதிகபட்சம்} mean? (\textquotedblleft{}At most; a maximum\textquotedblright{}.) Say it once more.
\end{tcolorbox}
\section[\texorpdfstring{\ta{முடிந்தவரை} (muṭintavarai)}{முடிந்தவரை (muṭintavarai)}]{\ta{முடிந்தவரை} (muṭintavarai) — as far as possible}

\label{lesson:TA-C554-mutintavarai}



\begin{quote}
Before the new one: say the Tamil for approximately, then the Tamil for at most.
\end{quote}

\subsection*{You'll want to know: \ta{முடிந்தவரை}}
\textbf{\ta{முடிந்தவரை}} — \emph{muṭintavarai} — \textquotedblleft{}as far as possible\textquotedblright{}.

\begin{quote}
\ta{முடிந்தவரை} \ta{சீக்கிரம்} \ta{கிளம்புவோம்}.
\end{quote}

\begin{quote}
\emph{muṭintavarai cīkkiram kiḷampuvōm.}
\end{quote}

\begin{quote}
We will set off as early as possible.
\end{quote}

Word for word \textquotedblleft{}up to what is possible\textquotedblright{}: \ta{முடிந்த} + \ta{வரை} (\emph{varai}).

\begin{cousinweb}
\textbf{\ta{முடிந்தவரை}} (\emph{muṭintavarai}) is a native phrase grown into one word: \ta{முடிந்த} (\emph{muṭinta}), \textquotedblleft{}that has been possible\textquotedblright{}, from \ta{முடி}, \textquotedblleft{}to end, to be accomplished\textquotedblright{} (DEDR 4922), and \ta{வரை} (\emph{varai}), \textquotedblleft{}up to\textquotedblright{}.
\end{cousinweb}

\begin{grammarlens}[title={What you've built}]
One more word for saying what happens if.
\end{grammarlens}

\subsection*{Your turn}
\begin{itemize}
\raggedright
  \item \emph{Say it:} \emph{muṭintavarai}
  \item \emph{Say it:} \emph{muṭintavarai}, once more
  \item \emph{Say it:} \emph{muṭintavarai cīkkiram kiḷampuvōm.} — We will set off as early as possible.
  \item \emph{Recall:} say the Tamil for approximately, then the Tamil for at most, then say \emph{muṭintavarai} again
\end{itemize}

\begin{tcolorbox}[breakable,colback=teal!4,colframe=teal!35!black,title={Before you move on}]
What does \ta{முடிந்தவரை} mean? (\textquotedblleft{}As far as possible\textquotedblright{}.) Say it once more.
\end{tcolorbox}
\section[\texorpdfstring{\ta{மலையேற்றம்} (malaiyēṟṟam)}{மலையேற்றம் (malaiyēṟṟam)}]{\ta{மலையேற்றம்} (malaiyēṟṟam) — trekking; hill climbing}

\label{lesson:TA-C554-malaiyerram}



\begin{quote}
Before the new one: say the Tamil for at most, then the Tamil for as far as possible.
\end{quote}

\subsection*{You'll want to know: \ta{மலையேற்றம்}}
\textbf{\ta{மலையேற்றம்}} — \emph{malaiyēṟṟam} — \textquotedblleft{}trekking; hill climbing\textquotedblright{}.

\begin{quote}
\ta{வானிலை} \ta{நன்றாக} \ta{இருந்தால்} \ta{சனிக்கிழமை} \ta{மலையேற்றம்} \ta{போவோம்}.
\end{quote}

\begin{quote}
\emph{vāṉilai naṉṟāka iruntāl caṉikkiḻamai malaiyēṟṟam pōvōm.}
\end{quote}

\begin{quote}
If the weather is good, we will go trekking on Saturday.
\end{quote}

\ta{மலை} (\emph{malai}) \textquotedblleft{}a hill\textquotedblright{} + \ta{ஏற்றம்} \textquotedblleft{}a climb\textquotedblright{}.

\begin{cousinweb}
\textbf{\ta{மலையேற்றம்}} (\emph{malaiyēṟṟam}) is a native compound: \ta{மலை} (\emph{malai}), \textquotedblleft{}a hill, a mountain\textquotedblright{} (DEDR 4742), and \ta{ஏற்றம்} (\emph{ēṟṟam}), \textquotedblleft{}a climb, an ascent\textquotedblright{}, from \ta{ஏறு}, \textquotedblleft{}to climb\textquotedblright{} (DEDR 916).
\end{cousinweb}

\begin{grammarlens}[title={What you've built}]
One more word for saying what happens if.
\end{grammarlens}

\subsection*{Your turn}
\begin{itemize}
\raggedright
  \item \emph{Say it:} \emph{malaiyēṟṟam}
  \item \emph{Say it:} \emph{malaiyēṟṟam}, once more
  \item \emph{Say it:} \emph{vāṉilai naṉṟāka iruntāl caṉikkiḻamai malaiyēṟṟam pōvōm.} — If the weather is good, we will go trekking on Saturday.
  \item \emph{Recall:} say the Tamil for at most, then the Tamil for as far as possible, then say \emph{malaiyēṟṟam} again
\end{itemize}

\begin{tcolorbox}[breakable,colback=teal!4,colframe=teal!35!black,title={Before you move on}]
What does \ta{மலையேற்றம்} mean? (\textquotedblleft{}Trekking; hill climbing\textquotedblright{}.) Say it once more.
\end{tcolorbox}
\section[\texorpdfstring{\ta{மீன்பிடி} (mīṉpiṭi)}{மீன்பிடி (mīṉpiṭi)}]{\ta{மீன்பிடி} (mīṉpiṭi) — to go fishing; to fish}

\label{lesson:TA-C554-minpiti}



\begin{quote}
Before the new one: say the Tamil for as far as possible, then the Tamil for trekking.
\end{quote}

\subsection*{You'll want to know: \ta{மீன்பிடி}}
\textbf{\ta{மீன்பிடி}} — \emph{mīṉpiṭi} — \textquotedblleft{}to go fishing; to fish\textquotedblright{}.

\begin{quote}
\ta{நேரம்} \ta{இருந்தால்} \ta{ஏரியில்} \ta{மீன்பிடிப்போம்}.
\end{quote}

\begin{quote}
\emph{nēram iruntāl ēriyil mīṉpiṭippōm.}
\end{quote}

\begin{quote}
If there is time, we will fish in the lake.
\end{quote}

Past (I …): \ta{மீன்பிடித்தேன்}, \emph{mīṉpiṭittēṉ}.

Word for word \textquotedblleft{}to catch fish\textquotedblright{}: \ta{மீன்} (\emph{mīṉ}) + \ta{பிடி} (\emph{piṭi}).

\begin{cousinweb}
\textbf{\ta{மீன்பிடி}} (\emph{mīṉpiṭi}) is a native compound: \ta{மீன்} (\emph{mīṉ}), \textquotedblleft{}a fish\textquotedblright{} (DEDR 4885), and \ta{பிடி} (\emph{piṭi}), \textquotedblleft{}to catch\textquotedblright{}.
\end{cousinweb}

\begin{grammarlens}[title={What you've built}]
One more word for saying what happens if.
\end{grammarlens}

\subsection*{Your turn}
\begin{itemize}
\raggedright
  \item \emph{Say it:} \emph{mīṉpiṭi}
  \item \emph{Say it:} \emph{mīṉpiṭi}, once more
  \item \emph{Say it:} \emph{nēram iruntāl ēriyil mīṉpiṭippōm.} — If there is time, we will fish in the lake.
  \item \emph{Recall:} say the Tamil for as far as possible, then the Tamil for trekking, then say \emph{mīṉpiṭi} again
\end{itemize}

\begin{tcolorbox}[breakable,colback=teal!4,colframe=teal!35!black,title={Before you move on}]
What does \ta{மீன்பிடி} mean? (\textquotedblleft{}To go fishing; to fish\textquotedblright{}.) Say it once more.
\end{tcolorbox}
\section[\texorpdfstring{\ta{பொம்மலாட்டம்} (pommalāṭṭam)}{பொம்மலாட்டம் (pommalāṭṭam)}]{\ta{பொம்மலாட்டம்} (pommalāṭṭam) — a puppet show}

\label{lesson:TA-C554-pommalattam}



\begin{quote}
Before the new one: say the Tamil for trekking, then the Tamil for to go fishing.
\end{quote}

\subsection*{You'll want to know: \ta{பொம்மலாட்டம்}}
\textbf{\ta{பொம்மலாட்டம்}} — \emph{pommalāṭṭam} — \textquotedblleft{}a puppet show\textquotedblright{}.

\begin{quote}
\ta{குழந்தைகள்} \ta{விரும்பினால்} \ta{பொம்மலாட்டம்} \ta{பார்க்கப்} \ta{போகலாம்}.
\end{quote}

\begin{quote}
\emph{kuḻantaikaḷ virumpiṉāl pommalāṭṭam pārkkap pōkalām.}
\end{quote}

\begin{quote}
If the children want to, we can go and see a puppet show.
\end{quote}

\ta{பொம்மை} (\emph{pommai}) \textquotedblleft{}a doll\textquotedblright{} + \ta{ஆட்டம்} \textquotedblleft{}a dance, a play\textquotedblright{}.

\begin{cousinweb}
\textbf{\ta{பொம்மலாட்டம்}} (\emph{pommalāṭṭam}) is a native compound: \ta{பொம்மை} (\emph{pommai}), \textquotedblleft{}a puppet, a doll\textquotedblright{} (DEDR 4530), and \ta{ஆட்டம்} (\emph{āṭṭam}), \textquotedblleft{}a dance, a play\textquotedblright{}.
\end{cousinweb}

\begin{grammarlens}[title={What you've built}]
One more word for saying what happens if.
\end{grammarlens}

\subsection*{Your turn}
\begin{itemize}
\raggedright
  \item \emph{Say it:} \emph{pommalāṭṭam}
  \item \emph{Say it:} \emph{pommalāṭṭam}, once more
  \item \emph{Say it:} \emph{kuḻantaikaḷ virumpiṉāl pommalāṭṭam pārkkap pōkalām.} — If the children want to, we can go and see a puppet show.
  \item \emph{Recall:} say the Tamil for trekking, then the Tamil for to go fishing, then say \emph{pommalāṭṭam} again
\end{itemize}

\begin{tcolorbox}[breakable,colback=teal!4,colframe=teal!35!black,title={Before you move on}]
What does \ta{பொம்மலாட்டம்} mean? (\textquotedblleft{}A puppet show\textquotedblright{}.) Say it once more.
\end{tcolorbox}
